\subsection{\red{Training-Time Collaborative Detection}}

\red{In conventional federated learning, clients retain local data while jointly building or updating a base model~\cite{federated_learning,Fedprox,ditto2021}. Recent NIDS systems follow this training-time pattern across detector designs~\cite{pourahmadi2023fedoe,mao2025fedinnid}. For example, Entente jointly trains graph-based detectors from distributed network logs~\cite{xu2026entente}. Federated tree methods likewise build their base models during collaboration~\cite{federboost,dpfedtree_ccs22}. Thus, these methods require local data while training the base detector. We instead freeze every detector before collaboration and learn only a personalized composer over its outputs.}

\red{Existing methods also adapt network-traffic classifier or detector training to specific deployment problems, such as noisy labels or network shift~\cite{qing2024rapier,xie2023rosetta,yang2026cold}. For example, CoLD filters noisy labels before training a detector~\cite{yang2026cold}. These methods modify the detector itself. We instead keep every detector frozen and learn a composer from limited current site data.}

\subsection{\red{Frozen Model Reuse and Personalization}}

\red{Some methods reuse pretrained models when their original training data are unavailable. For example, Old Dogs learns one target classifier from pretrained source networks without using source data~\cite{old_dogs2019}. In a multiparty setting, HMR exchanges selected examples and calibrates the reused models during collaboration~\cite{wu2019hmr}. Likewise, Ma et al. keep local models frozen but send per-example representations of aligned samples to a central server~\cite{ma2023consensus}. We instead keep all detectors frozen, keep per-example outputs local, and make adoption site-specific.}

\red{Other methods first train local classifiers from current data and then learn a second-stage model. F-BIDS and Fed-MStacking follow this pattern~\cite{aouedi2023fbids,qiu2024fedmstacking}. Meanwhile, personalized federated learning updates base models during collaboration~\cite{fedfomo2021,ditto2021}. We instead start from detectors frozen before collaboration and train only the composer.}

\subsection{\red{Protected Collaboration and Local Adoption}}

\red{The two stages raise different questions. During candidate construction, keeping current site data local does not protect the composer gradient sent by each site. We therefore use secure aggregation so that the coordinator receives only the aggregate gradient and no individual gradient in plaintext when sites do not collude with it~\cite{bonawitz2017secureaggregation}. By comparison, FedDef perturbs training inputs to reduce gradient-reconstruction leakage in federated NIDSs~\cite{chen2023feddef}. Moreover, differential privacy limits how much the output distribution can change between neighboring datasets~\cite{dwork2006calibrating}. However, our protocol provides neither differential privacy nor protection against inference from the aggregates.}

\red{Deployment authorization begins only after construction ends. When old training data are unavailable or do not match the target distribution, Active Comparison uses new labels to compare a fixed baseline and challenger~\cite{sawade2012activecomparison}. Separately, Learn Then Test evaluates prespecified risk hypotheses on independent calibration data~\cite{angelopoulos2025ltt}. We likewise freeze the candidate before using audit data, but couple this decision to the preceding example-local construction stage and bound excess risk relative to each site's audited incumbent. Existing protection and comparison mechanisms therefore address the two stages separately rather than the complete post-training deployment contract.}
